\section{Chương~\ref{sec:recognition}. Nhận dạng}

Tốc độ nhận dạng, cấu tạo của các tầng đầu, việc đọc tuyến tính câu trả lời và việc học từ sự liên tục của thời gian đã được đo ở người và khỉ; việc quy cả chuỗi về hai phép toán và việc học từ video đã được kiểm tra trên các mô hình của sách, còn cách giải thích các ảo giác đi từ có cơ sở vững đến còn tranh cãi.

{\footnotesize
\setlength{\tabcolsep}{3pt}\renewcommand{\arraystretch}{1.15}
\begin{longtable}{>{\raggedright}p{0.5cm}>{\raggedright}p{4.0cm}>{\raggedright}p{5.6cm}>{\raggedright}p{2.3cm}>{\raggedright\arraybackslash}p{1.7cm}}
\toprule
TT & Khẳng định & Thí nghiệm và điều đã đo & Nguồn & Trạng thái \\
\midrule\endfirsthead
\toprule
TT & Khẳng định & Thí nghiệm và điều đã đo & Nguồn & Trạng thái \\
\midrule\endhead
1 & So sánh với mẫu theo điểm ảnh khi có dịch chuyển đoán đúng không hơn đồng xu; cặp tầng~\eqref{eq:scstage} nhận ra hình ở bất kỳ vị trí nào & \texttt{models/\allowbreak recognition\_models.py}, ``võng mạc'' $24$ điểm, $4000$ lần thử: so mẫu cho $0.49$, $0.51$, $0.52$ với tỷ lệ điểm bị lật $0$, $0.1$, $0.2$; cặp $S$ và $C$ cho $1.00$, $0.86$, $0.70$ & suy ra trong chương & mô hình \\
2 & Con vật trong ảnh được nhận ra sau $\sim150$\,ms, một lượt qua chừng mười tầng, mỗi tầng $10$--$20$\,ms & Người, nhiệm vụ ``có con vật hay không'' trên ảnh hiện $20$\,ms: điện thế gợi của hai câu trả lời tách nhau sau $\sim150$\,ms. Khỉ: thời điểm đáp ứng đầu tiên tăng dần từ tầng này sang tầng khác (V1 sớm nhất, rồi V2 và V4). Số tầng và ``không hoặc một xung mỗi tầng'' được suy ra từ các con số này & \cite{Thorpe1996,Schmolesky1998} & đã đo \\
3 & Tầng $S$ là VÀ theo bộ phận, tầng $C$ là giá trị lớn nhất theo vị trí (mệnh đề~\ref{prop:conveyor}) & Mèo, vỏ não thị giác sơ cấp: tế bào đơn giản đáp ứng với cạnh có hướng nhất định ở vị trí nhất định, tế bào phức tạp đáp ứng với cùng hướng đó trong vùng rộng hơn. Ghi nội bào tế bào phức tạp: ở nhiều tế bào, đáp ứng với hai vạch gần bằng đáp ứng lớn hơn trong hai, trung bình gần nhất với phép max. Giá trị lớn nhất qua ức chế lẫn nhau: Mệnh đề~\ref{prop:wta}; phép chia cho hoạt động chung: tổng quan & \cite{HubelWiesel1962,Lampl2004,Carandini2012} & đã đo \\
4 & Càng lên cao trong chuỗi, tổ hợp càng lớn và phức tạp, câu trả lời càng ít phụ thuộc vị trí & Khỉ, đơn giản hóa kích thích tới ``đặc trưng tới hạn'': tế bào V4 và vỏ não thái dương dưới phía sau đáp ứng với tổ hợp phức tạp của hình dạng, màu sắc và kết cấu, với trường tiếp nhận nhỏ; ở vỏ não thái dương dưới phía trước, trường lớn hơn. Sự xen kẽ $S$ và $C$ trong mô hình tái hiện bức tranh này & \cite{KobatakeTanaka1994,Riesenhuber1999} & đã đo \\
5 & Mạng tích chập lặp lại sự xen kẽ này và dự đoán tốt nhất đáp ứng của các tầng trên; đồ vật được đọc tuyến tính từ tần số của một nhóm nơ-ron & Mạng được huấn luyện để nhận dạng, không khớp với não: lớp trên cùng dự đoán đáp ứng của nơ-ron vỏ não thái dương dưới của khỉ, các lớp giữa dự đoán đáp ứng của V4. Bộ phân loại tuyến tính dựa trên hoạt động của $\sim100$ tế bào chọn ngẫu nhiên, trong cửa sổ từ $12.5$\,ms, nhận ra đồ vật và phạm trù ở các vị trí và kích thước khác nhau & \cite{Fukushima1980,Yamins2014,Hung2005} & đã đo \\
6 & Quy tắc Hebb có vết~\eqref{eq:traceinvariance} dạy tính bất biến từ video không nhãn nếu vết không ngắn hơn $\sim20$\,ms & Ý tưởng đặc trưng chậm và quy tắc có vết là lý thuyết. \texttt{models/\allowbreak recognition\_models.py}, hai nơ-ron $C$, $16$ đầu vào, $10$ lần chạy: nghỉ giữa các đồ vật $0.3$\,s. Với $\tau_{\text{tr}}=5$\,ms, trùng với hình trung bình $0.52$, với $10$\,ms là $0.56$; với $20$, $50$ và $200$\,ms là $1.00$ ở mọi lần chạy (với $30$\,ms, một trong mười lần chạy sai) & \cite{Foldiak1991,WiskottSejnowski2002} & mô hình \\
7 & Trong não, tính bất biến được học từ sự liên tục của thời gian & Khỉ: đồ vật bị tráo kín đáo trong lúc mắt nhảy tới một vị trí nhất định của trường nhìn; tính bất biến vị trí của nơ-ron vỏ não thái dương dưới thay đổi theo sự đánh tráo, có ý nghĩa thống kê chỉ sau một giờ. Việc đây là quy tắc chính của học thị giác là giả thuyết & \cite{LiDiCarlo2008} & đã đo \\
8 & Chuyển động: bộ phát hiện hướng, rồi cùng cặp VÀ/HOẶC trên những vùng lớn & Bộ phát hiện hướng trong võng mạc thỏ; ở vùng MT của khỉ, cả $110$ nơ-ron được ghi đều chọn lọc theo hướng, đáp ứng với chuyển động mạnh hơn ở V1 & \cite{Barlow1965,Albright1984} & đã đo \\
9 & Các tầng trên gửi xuống một dự đoán, đi lên là sai số & Mô hình giải thích các hiệu ứng ngoài trường cổ điển ở vỏ não thị giác sơ cấp; một số quan sát riêng lẻ phù hợp (tổng quan), như một cấu tạo chung của dây chuyền thì chưa được chứng minh & \cite{RaoBallard1999,Keller2018} & giả thuyết \\
10 & Ảnh khó cần phản hồi và vài vòng & Khỉ: với ảnh ``khó'', quyết định ở vỏ não thái dương dưới hình thành muộn hơn $\sim30$\,ms; các đáp ứng muộn này được mạng đi thẳng dự đoán kém, còn mạng có phản hồi hoặc mạng sâu hơn dự đoán tốt hơn & \cite{Kar2019} & đã đo \\
11 & Sửa điểm ảnh không nhận ra được đánh lừa mạng; thị giác người bền vững hơn nhiều nhưng không miễn nhiễm & Trên mạng: một nhiễu loạn nhỏ được chọn đặc biệt làm đổi câu trả lời; ví dụ ``gấu trúc $\to$ vượn'' lấy từ công trình thứ hai. Ở người, khi hiện ngắn, các ảnh sửa có khả năng chuyển tốt giữa các mạng làm lệch lựa chọn & \cite{Szegedy2014,Goodfellow2015,Elsayed2018} & đã đo \\
12 & Ảo giác kỳ vọng: đầu vào không đáng tin nhường cho kỳ vọng; vật nhạt chuyển động ``chậm hơn'', chiếc váy được thấy khác nhau (mục~\ref{sec:illusions}) & Cách tử có độ tương phản thấp hơn trông như chuyển động chậm hơn. Khảo sát $1401$ người: chiếc váy được thấy theo ba phương án ổn định, gợi ý về ánh sáng làm đổi màu; ở $\sim13\,000$ người quan sát, giả định về ánh sáng quyết định. Mô hình Bayes với kỳ vọng về chuyển động chậm giải thích một loạt ảo giác chuyển động & \cite{Thompson1982,LaferSousa2015,Wallisch2017,Weiss2002,Purves2011} & giả thuyết \\
13 & Điều chỉnh khuếch đại: thác nước, độ nghiêng và độ tương phản; lưới Hermann không phải ức chế láng giềng & Bộ phát hiện hướng ở võng mạc thỏ giảm tần số khi chuyển động kéo dài, và sau khi dừng thì tụt xuống dưới mức nền. Ảo giác độ nghiêng được tái hiện bằng mô hình chia cho hoạt động của láng giềng. Những biến đổi lưới không làm đổi đáp ứng của nơ-ron đầu ra võng mạc lại làm mất các đốm: cách giải thích bằng ức chế láng giềng bị bác bỏ & \cite{BarlowHill1963,Schwartz2009,SchillerCarvey2005} & đã đo \\
14 & Bù độ trễ: vật đang chuyển động có vẻ nằm trước điểm lóe & Ảo giác đã được đo. Tâm vật lý học mâu thuẫn với cả việc ngoại suy chuyển động lẫn chênh lệch độ trễ; một cách giải thích hồi tố đã được đề xuất, dựa trên các sự kiện $\sim80$\,ms sau điểm lóe. Tranh luận chưa ngã ngũ & \cite{Nijhawan1994,EaglemanSejnowski2000} & giả thuyết \\
15 & ``Rắn'' đứng yên lại quay: chênh lệch độ trễ~\eqref{eq:latency} được bộ phát hiện hướng đọc thành chuyển động & Ảo giác vẫn xảy ra khi không có màu; từng cặp phần tử kề nhau tự tạo ra nó; nơ-ron vỏ não khỉ chọn lọc theo hướng cho đáp ứng có hướng với các cặp đứng yên này. Ở người, chuyển động quay được kích hoạt bởi vi giật mắt và cái chớp mắt & \cite{Conway2005,OteroMillan2012} & đã đo \\
16 & Ảnh hai nghĩa: ức chế lẫn nhau, sự mỏi và nhiễu; sự chuyển chủ yếu do nhiễu gây ra & Mô hình các nhóm nơ-ron cạnh tranh: sự chuyển trong mô hình do \emph{nhiễu} gây ra, không có nhiễu thì không chuyển; mô hình giải thích việc tần số chuyển tăng theo cường độ kích thích. Sự mỏi ở đây đóng vai trò nhỏ hơn & \cite{MorenoBote2007} & giả thuyết \\
17 & Một phần ảo giác là hệ quả của bài toán mà cỗ máy học & Mạng được huấn luyện để dự đoán khung hình tiếp theo của video ``thấy'' các vòng tròn tĩnh quay và không thấy điều đó trên ảnh đối chứng; mạng khử nhiễu và làm nét ảnh mắc các ảo giác màu sắc và độ sáng như con người & \cite{Watanabe2018,GomezVilla2020} & đã đo \\
\bottomrule
\end{longtable}}
